%==============================================================================
%  CAPÍTULO X
%==============================================================================
\chapter{La paradoja del \texorpdfstring{\emph{discharge}}{discharge} y la exclusión del
Crédito con Aval del Estado}
\label{cap:discharge-cae}

\fichaCapitulo{Incompatibilidad de las deudas de estudio con el perdón concursal.}{Tensión
entre universalidad y especialidad, y evolución de la jurisprudencia de la Corte Suprema.}

%==============================================================================
\bloqueTeorico
%==============================================================================

\subsection{El conflicto de principios}

Colisión entre el beneficio general del \discharge{} (\art{255}) ---dirigido a liberar de
cargas perpetuas al individuo y reinsertarlo productivamente--- y la protección del
interés fiscal comprometido en el financiamiento de la educación superior (\Lcae{} y
\Lfondosolidario{}).

\subsection{La ausencia de un catálogo legal de deudas excluidas}
\pendiente{Desarrollar: cómo la falta de codificación expresa desplaza la solución hacia
la interpretación judicial.}

%==============================================================================
\bloqueNormativo
%==============================================================================

\subsection{Universalidad subjetiva y efecto de remisión de la resolución de término}
\pendiente{Desarrollar.}

\subsection{El principio de especialidad}

Aplicación del \art{8} en relación con las reglas del \cc{} sobre prevalencia de la ley
especial.\verificar{Confirmar el contenido del art. 8 y su idoneidad para fundar la
especialidad invocada.}

\subsection{El régimen del Crédito con Aval del Estado}

Naturaleza del préstamo, garantías fiscales ejecutables por la \tgr{} y cobro contingente
al ingreso, conforme a la \Lcae{}.

\subsection{El Crédito Solidario Universitario}

Régimen de la \Lfondosolidario{}.

%==============================================================================
\bloquePractico
%==============================================================================

\subsection{Secuencia procesal del deudor con deudas educativas}

\begin{flujograma}{Suerte del CAE en el procedimiento de liquidación voluntaria}{cap10-cae}
  \node[hito] (sol) {EL DEUDOR PRESENTA SOLICITUD DE LIQUIDACIÓN VOLUNTARIA\\[2pt]
    \normalfont Debe declarar la deuda CAE en la nómina de pasivos};
  \node[nodo, below=of sol] (excl)
    {LA INSTITUCIÓN BANCARIA O LA \tgr{} SOLICITA LA EXCLUSIÓN DEL CRÉDITO};

  \coordinate (split) at ($(excl.south)+(0,-6mm)$);
  \node[adverso, below left=12mm and 4mm of excl.south] (acoge)
    {EL TRIBUNAL ACOGE LA EXCLUSIÓN\\[2pt]
     \normalfont La deuda sobrevive a la liquidación};
  \node[favorable, below right=12mm and 4mm of excl.south] (rechaza)
    {EL TRIBUNAL RECHAZA LA EXCLUSIÓN\\[2pt]
     \normalfont La deuda ingresa al concurso y se sujeta al \art{255}};

  \node[nodo, text width=40mm, below=9mm of acoge] (ape)
    {Recurso de apelación y, en su caso, acción de protección};
  \node[nodo, text width=40mm, below=9mm of rechaza] (cas)
    {Recurso de casación del acreedor o de la \tgr{}};

  \draw[flecha] (sol) -- (excl);
  \draw[flecha] (excl.south) -- (split) -| (acoge.north);
  \draw[flecha] (excl.south) -- (split) -| (rechaza.north);
  \draw[flecha] (acoge) -- (ape);
  \draw[flecha] (rechaza) -- (cas);
\end{flujograma}

\subsection{Estrategias de defensa del deudor}
\pendiente{Desarrollar según el estado actual de la jurisprudencia.}

%==============================================================================
\bloqueJurisprudencial
%==============================================================================

\subsection{Evolución cronológica de la jurisprudencia de la Corte Suprema}

\subsubsection{La postura clásica de la Sala Civil}

Fallo \emph{Jamarne con Salazar}, \rol{4656-2017}, que ordenó la exclusión del CAE
aplicando la primacía de la ley especial de educación superior por sobre el régimen
general concursal.\verificar{Confirmar carátula, rol, sala y fecha.}

\subsubsection{El giro de la Tercera Sala}

Sentencia de protección \emph{Mancilla con Tesorería General de la República},
\rol{59.567-2020}, que amparó al deudor sosteniendo que el CAE no constituye un estatuto
concursal especial de insolvencia y que, por aplicación de la regla de la \lat{lex
posterior}, la \LIR{} debe extinguir la totalidad de los pasivos.\verificar{Confirmar
carátula, rol, sala y fecha.}

\subsubsection{El retorno a la exclusión}

Sentencias \rol{3557-2025} y \rol{47.641-2024} de la Sala Civil, que confirmaron la
improcedencia de la descarga del CAE por vía concursal, fijando la doctrina de que la
deuda estudiantil sobrevive íntegramente al procedimiento. Debe incorporarse el voto
disidente que limita la facultad de invocar la especialidad únicamente a la \tgr{} y no a
los bancos de la plaza.\verificar{Confirmar roles, fechas y la individualización del
ministro disidente antes de publicar.}

\begin{alerta}[Estado de la cuestión]
La materia es de criterio jurisprudencial cambiante y con sentencias muy recientes. Toda
afirmación de este capítulo debe contrastarse con la jurisprudencia vigente a la fecha de
cierre de la edición.
\end{alerta}
